Dependence polo\"{\i}dale du champ electrique Ohmique

Le fonctionnement de base du tokamak repose sur le principe du transformateur,
l'anneau toro\"{\i}dal forme par le plasma constituant le circuit
secondaire. Dans ce regime, le courant total $I_{p}$ circulant dans
le plasma est induit par la variation temporelle du flux polo\"{\i}dal
$\frac{\partial\psi}{\partial t}$ produite par le circuit primaire.
La densite de courant $\overrightarrow{j}$ circulant dans le plasma
resulte de la loi d'Ohm, $\overrightarrow{j}=\overline{\overline{\sigma}}_{\Omega}\overrightarrow{E},$
ou $\overline{\overline{\sigma}}_{\Omega}$ est la conductivite electrique
du plasma prenant en compte les interactions Coulombiennes entre ses
constituants, les electrons etant fortement magnetises. La relation
entre $\frac{\partial\psi}{\partial t}$ et $\overrightarrow{E}$
est obtenue a partir de la loi Faraday, \[
\frac{\partial\overrightarrow{B}}{\partial t}=-\overrightarrow{\nabla}\times\overrightarrow{E}\]
 En projetant cette expression sur le gradient du flux polo\"{\i}dal,
on obtient, \[
\overrightarrow{\nabla}\psi\cdot\frac{\partial\overrightarrow{B}}{\partial t}=-\overrightarrow{\nabla}\psi\cdot\left(\overrightarrow{\nabla}\times\overrightarrow{E}\right)\]
 et en utilisant le fait que le champ magnetique est contenu dans
les surfaces de flux, \[
\overrightarrow{B}\cdot\overrightarrow{\nabla}\psi=0\rightarrow\overrightarrow{\nabla}\psi\cdot\frac{\partial\overrightarrow{B}}{\partial t}=-\overrightarrow{B}\cdot\overrightarrow{\nabla}\frac{\partial\psi}{\partial t}\]
 ce qui conduit a \[
\overrightarrow{B}\cdot\overrightarrow{\nabla}\frac{\partial\psi}{\partial t}=\overrightarrow{\nabla}\psi\cdot\left(\overrightarrow{\nabla}\times\overrightarrow{E}\right)\]


Sachant que \[
\overrightarrow{\nabla}\times\overrightarrow{E}=\overrightarrow{\nabla}\times E_{\phi}\overrightarrow{e}_{\phi}\]
 en raison de l'axisymetrie, on obtient\[
\overrightarrow{\nabla}\psi\cdot\left(\overrightarrow{\nabla}\times\overrightarrow{E}\right)=\overrightarrow{\nabla}\psi\cdot\left(\overrightarrow{\nabla}\times E_{\phi}\overrightarrow{e}_{\phi}\right)=-\left(\overrightarrow{\nabla}\psi\times\overrightarrow{e}_{\phi}\right)\cdot\overrightarrow{\nabla}E_{\phi}\]
en utilisant les relations vectorielles suivantes: $\overrightarrow{\nabla}\cdot\left(\overrightarrow{A}\times\overrightarrow{B}\right)=\overrightarrow{B}\cdot\left(\overrightarrow{\nabla}\times\overrightarrow{A}\right)-\overrightarrow{A}\cdot\left(\overrightarrow{\nabla}\times\overrightarrow{B}\right)$
et $\nabla\times\left(f\overrightarrow{A}\right)=f\overrightarrow{\nabla}\times\overrightarrow{A}+\overrightarrow{\nabla}f\times\overrightarrow{A}$

D'apres l'expression du champ poloidal en fonction de $\psi,$ on
trouve donc \[
\overrightarrow{\nabla}\psi\cdot\left(\overrightarrow{\nabla}\times\overrightarrow{E}\right)=-R\overrightarrow{B}_{p}\cdot\overrightarrow{\nabla}E_{\phi}\]
 Comme de plus $\overrightarrow{B}_{\phi}\cdot\overrightarrow{\nabla}E_{\phi}=0,$
il est toujours possible d'ajouter ce terme au second membre, afin
de faire apparaitre le champ magnetique total $\overrightarrow{B}=$
$\overrightarrow{B}_{\phi}+\overrightarrow{B}_{p}.$ On en deduit
donc \[
\overrightarrow{\nabla}\psi\cdot\left(\overrightarrow{\nabla}\times\overrightarrow{E}\right)=-R\overrightarrow{B}\cdot\overrightarrow{\nabla}E_{\phi}=-\overrightarrow{B}\cdot\overrightarrow{\nabla}RE_{\phi}\]


Au final, \[
\overrightarrow{B}\cdot\overrightarrow{\nabla}\frac{\partial\psi}{\partial t}=-\overrightarrow{B}\cdot\overrightarrow{\nabla}RE_{\phi}\]
 soit \[
\overrightarrow{B}\cdot\overrightarrow{\nabla}\left(\frac{\partial\psi}{\partial t}+RE_{\phi}\right)=0\]
 quelle que soit la valeur de $\overrightarrow{B},$ ce qui indique
que \[
\frac{\partial\psi}{\partial t}=-RE_{\phi}\]


En integrant cette relation, en supposant que $E_{\phi}$ est une
fonction independante du temps, on trouve que $\psi\left(t\right)=-RE_{\phi}t+\psi\left(t=0\right).$
Sachant que $\overrightarrow{B}_{p}$ depend de $\overrightarrow{\nabla}\psi,$
on en deduit immediatement que $\frac{\partial\psi}{\partial R}$
$=-t\frac{\partial RE_{\phi}}{\partial R}+\frac{\partial\psi\left(t=0\right)}{\partial R}$
mais egalement $\frac{\partial\psi}{\partial Z}$ $=-t\frac{\partial RE_{\phi}}{\partial Z}+\frac{\partial\psi\left(t=0\right)}{\partial Z}.$
Le champ polo\"{\i}dal etant constant, puisqu'il est une fonction
de $\overrightarrow{j}=\overline{\overline{\sigma}}_{\Omega}\overrightarrow{E}$
selon la loi d'Ohm, cela implique que $RE_{\phi}$ est une constante
independante de $R$ et de $Z$, sinon $\overrightarrow{B}_{p}$ evoluerait
aussi en fonction du temps. On a ainsi \[
R\left(r,\theta\right)E_{\phi}\left(r,\theta\right)=R\left(r,\theta=0\right)E_{\phi}\left(r,\theta=0\right)\]
 soit \[
E_{\phi}\left(r,\theta\right)=\Psi\left(r,\theta\right)E_{\phi}^{(0)}\left(r\right)\]


Expression conservative du flux lie au champ electrique constant

En posant $\mathcal{E}(f_{0})=\overrightarrow{\nabla}\cdot\overrightarrow{S}^{E},$on
trouve immediatement que dans des coordonnees cylindriques $\left(p_{\Vert},p_{\bot}\right),$
le flux induit par le champ electrique longitudinal $E_{\Vert}$ dans
l'espace des vitesses vaut \[
\overrightarrow{S}^{E}=q_{e}E_{_{\Vert}}f_{0}\overrightarrow{b}\]
 soit \[
\overrightarrow{S}^{E}\left(cyl.\right)=\left(\begin{array}{l}
S_{p_{_{\Vert}}}^{E}\\
S_{p_{\bot}}^{E}\end{array}\right)=\left(\begin{array}{l}
q_{e}E_{\Vert}f_{0}\\
0\end{array}\right)\]
 ou $q_{e}$ est la charge electrique de l'electron, et $\overrightarrow{b}=$
$\overrightarrow{B}/\left|\left|\overrightarrow{B}\right|\right|$le
vecteur unite le long de la ligne de champ.

Pour calculer le flux dans en coordonnees spheriques $\left(p,\xi\right),$
ou $\xi=p_{_{\Vert}}/p,$ et $p_{\bot}/p=\sqrt{1-\xi^{2}},$ on effectue
la transformation $\overrightarrow{S}^{E}\left(sph.\right)=\mathcal{R}^{-1}$
$\overrightarrow{S}^{E}\left(cyl.\right),$ ou $\mathcal{R}$ est
la rotation permettant de passer du repere cylindrique au spherique
soit \[
\mathcal{R}=\left(\begin{array}{ll}
\xi & -\sqrt{1-\xi^{2}}\\
\sqrt{1-\xi^{2}} & \xi\end{array}\right)\]


Utilisant la propriete usuelle des rotations, $\mathcal{R}^{-1}=$
$^{t}\mathcal{R}$ on en deduit que \[
\left\{ \begin{array}{l}
S_{p}^{E}=\xi S_{p_{\Vert}}^{E}\\
S_{\xi}^{E}=-\sqrt{1-\xi^{2}}S_{p_{\Vert}}^{E}\end{array}\right.\]
 et naturellement \begin{eqnarray*}
F_{p}^{\text{E}} & = & q_{e}\xi E_{\Vert}\\
F_{\xi}^{\text{E}} & = & -q_{e}\sqrt{1-\xi^{2}}E_{\Vert}\end{eqnarray*}


A partir de ces expressions il est facile de determiner la moyenne
des coefficients sur la periode de rebond. On obtient ainsi

\begin{eqnarray*}
F_{p}^{\text{E}(0)} & = & q_{e}\left\{ \xi E_{\Vert}\right\} \\
F_{\xi}^{\text{E}(0)} & = & q_{e}\sigma\left\{ -\sqrt{1-\xi^{2}}\frac{\sigma\xi}{\sqrt{\Psi}\xi_{0}}E_{\Vert}\right\} \end{eqnarray*}
 Comme la dependence poloidale du champ electrique est la suivante
sur une surface de flux $r$ donnee \[
E_{\Vert}(\theta)=\Psi(\theta)E_{\Vert,0}\]
 on en deduit que

\begin{eqnarray*}
F_{p}^{\text{E}(0)} & = & \xi_{0}\sigma q_{e}E_{\Vert,0}\left\{ \Psi\frac{\sigma\xi}{\xi_{0}}\right\} \\
F_{\xi}^{\text{E}(0)} & = & -\sqrt{1-\xi_{0}^{2}}\sigma q_{e}E_{\Vert,0}\left\{ \Psi\frac{\sigma\xi}{\xi_{0}}\right\} \end{eqnarray*}


Le coefficient $\left\{ \Psi\xi/\xi_{0}\right\} $ est defini selon
la forme usuelle \[
\left\{ \Psi\frac{\sigma\xi}{\xi_{0}}\right\} \equiv\frac{\sigma s^{\ast}}{\lambda\left(\xi_{0}\right)}\]
 et on obtient finalement \begin{eqnarray*}
F_{p}^{\text{E}(0)} & = & \xi_{0}\frac{s^{\ast}}{\lambda\left(\xi_{0}\right)}q_{e}E_{\Vert,0}\\
F_{\xi}^{\text{E}(0)} & = & -\sqrt{1-\xi_{0}^{2}}\frac{s^{\ast}}{\lambda\left(\xi_{0}\right)}q_{e}E_{\Vert,0}\end{eqnarray*}
 Sous forme discrete, les coefficients lies au champ electrique constant
qui interviennent dans la matrice $\overline{\overline{M}}^{\left(0\right)}$
sont \[
\overline{\overline{D}}_{p}^{\text{E}(0)}\rightarrow\left\{ \begin{array}{l}
D_{pp,i+1,j+1/2}^{\text{E}\left(0\right)}=0\\
D_{pp,i,j+1/2}^{\text{E}\left(0\right)}=0\\
D_{p\xi,i+1,j+1/2}^{\text{E}\left(0\right)}=0\\
D_{\xi p,i+1/2,j+1}^{\text{E}\left(0\right)}=0\\
D_{p\xi,i+1/2,j+1/2}^{\text{E}\left(0\right)}=0\\
D_{\xi p,i+1/2,j+1/2}^{\text{E}\left(0\right)}=0\\
D_{p\xi,i,j+1/2}^{\text{E}\left(0\right)}=0\\
D_{\xi p,i+1/2,j}^{\text{E}\left(0\right)}=0\\
D_{\xi\xi,i+1/2,j+1}^{\text{E}\left(0\right)}=0\\
D_{\xi\xi,i+1/2,j}^{\text{E}\left(0\right)}=0\end{array}\right.\]
 et \[
\overline{F}_{p}^{\text{E}(0)}\rightarrow\left\{ \begin{array}{l}
F_{p,i+1,j+1/2}^{\text{E}(0)}=\xi_{0j+1/2}\frac{s_{j+1/2}^{\ast}}{\lambda\left(\xi_{0j+1/2}\right)}(q_{e}E_{\Vert,0})\\
F_{p,i,j+1/2}^{\text{E}(0)}=\xi_{0j+1/2}\frac{s_{j+1/2}^{\ast}}{\lambda\left(\xi_{0j+1/2}\right)}(q_{e}E_{\Vert,0})\\
F_{\xi,i+1/2,j+1}^{\text{E}(0)}=-\sqrt{1-\xi_{0j+1}^{2}}\frac{s_{j+1}^{\ast}}{\lambda\left(\xi_{0j+1}\right)}(q_{e}E_{\Vert,0})\\
F_{\xi,i+1/2,j}^{\text{E}(0)}=-\sqrt{1-\xi_{0j}^{2}}\frac{s_{j}^{\ast}}{\lambda\left(\xi_{0j}\right)}(q_{e}E_{\Vert,0})\end{array}\right.\]

